% =============================================================================
% 03_results.tex -- Results section
% Figures 1--6 and 9: code/calc-01.py .. calc-06.py, calc-09.py -> manuscript/figures/figure_1..6, figure_9.
% Figs.~7--8 live in Discussion. Re-run "make figures" or "scripts\\build.ps1 -Figures" to refresh them.
% =============================================================================

\section{Results}
\label{sec:results}

\subsection{Key Analytical Results}
\label{subsec:analytical-results}

The covariant phase-locked Klein--Gordon (PLKG) equations \eqref{eq:14} and \eqref{eq:16}
yield the following key results:
\begin{itemize}
    \item Effective synchronized phase sector: $w_{\text{sync}} \approx 1$
          when coherent phase-locking admits an effective barotropic scaling $\rho\propto a^{-6}$, i.e.\ $w=1$ with negligible energy exchange with other sectors.
          This stiff regime requires the PLKG stress tensor to be dominated by coherent phase-kinetic contributions.
          The bare Kuramoto interaction potential in $\mathcal{L}$, viewed in isolation, is vacuum-like with $w_{V_{\text{int}}}=-1$, not $w_{\text{sync}}\approx 1$.
    \item Effective synchronization length:
          $\lambda_{\text{sync}} \sim \dfrac{1}{\sqrt{K}}$
    \item Modified dispersion relation (decoupled amplitude--phase limit; Appendix~\ref{app:stability}, Eqs.~\eqref{eq:dispersion} and~\eqref{eq:meff}):
      \begin{equation*}
      \omega^2 + 3iH\omega - \frac{k^2}{a^2} - m_{\text{eff}}^2 = 0,
      \end{equation*}
      with $m_{\text{eff}}^2 = m^2 - K - \dot{\bar{\theta}}^2$.
\end{itemize}

\subsection{Numerical Illustration of Synchronization}
\label{subsec:numerical}

To illustrate the qualitative behaviour predicted by Eqs.~\eqref{eq:14} and~\eqref{eq:16}, the simplified point-dynamics version of the PLKG system (Hubble friction omitted,
spatial gradients suppressed) is integrated for $N=5$ oscillators with masses
$m_j \in [1.0,\,1.5]$ and coupling $K=2.5$.
There is no spatial domain in this truncation, so no spatial boundary data are imposed.
The only side data are autonomous initial conditions at $t=0$.
Initial values $\theta_j(0)$ are drawn independently and uniformly on $[-\pi,\pi]$ and $\rho_j(0)$ independently and uniformly on $[0.8,\,1.2]$, using a fixed pseudorandom seed shared across runs so Figs.~\ref{fig:kkg-sync}--\ref{fig:kkg-K-scan} are reproducible and directly comparable.
The driver script is \texttt{code/calc-01.py} (Appendix~\ref{app:reproducibility}).

\begin{figure}[t]
    \centering
    \includegraphics[width=\columnwidth]{figures/figure_1.pdf}
    \caption{Local PLKG dynamics for $N=5$ coupled scalar oscillators with
    coupling $K=2.5$ in the homogeneous ODE truncation ($H=0$, $\nabla\to 0$), so no spatial boundary-value problem is specified.
    Initial data at $t=0$: $\theta_j(0)\sim\mathcal{U}(-\pi,\pi)$ and $\rho_j(0)\sim\mathcal{U}(0.8,1.2)$ i.i.d., fixed seed.
    (a) Phase observables $\sin\theta_j(t)$; (b) amplitudes $\rho_j(t)$; (c) Kuramoto order parameter $R(t)=|N^{-1}\sum_j e^{i\theta_j}|$.
    The order parameter approaches unity, demonstrating mean-field phase locking.
    Generated by \texttt{code/calc-01.py}.}
    \label{fig:kkg-sync}
\end{figure}

The same simplified dynamics makes the approach of $R(t)$ toward coherence sensitive to the coupling: larger $K$ tends to shorten the transient toward high $R$, whereas small $K$ leaves the system nearer the disordered manifold for longer.
Figure~\ref{fig:kkg-order-K} overlays $R(t)$ for several choices of $K$ with the same truncation, masses, and the same initial draw $(\theta_j(0),\rho_j(0))$ as in Fig.~\ref{fig:kkg-sync} (\texttt{code/calc-02.py}).

\begin{figure*}[t]
    \centering
    \includegraphics[width=\textwidth]{figures/figure_2.pdf}
    \caption{Kuramoto order parameter $R(t)=|N^{-1}\sum_j e^{i\theta_j}|$ for the same homogeneous ODE setup as Fig.~\ref{fig:kkg-sync} ($H=0$, $\nabla\to 0$; no spatial boundaries), with $N=5$ and $m_j\in[1.0,1.5]$.
    Each curve uses the same initial draw $\theta_j(0)\sim\mathcal{U}(-\pi,\pi)$, $\rho_j(0)\sim\mathcal{U}(0.8,1.2)$ and varies only $K$.
    Larger $K$ drives faster macroscopic phase alignment.
    Generated by \texttt{code/calc-02.py}.}
    \label{fig:kkg-order-K}
\end{figure*}

Aggregating many couplings at fixed integration time makes the crossover in macroscopic coherence clearer.
Figure~\ref{fig:kkg-K-scan} scans $K$ on a dense grid from $0.05$ to $12$ with the same homogeneous truncation, masses, initial draw, and seed as Figs.~\ref{fig:kkg-sync} and \ref{fig:kkg-order-K}.
A fixed horizon $t_{\mathrm{f}}=45$ is held, matching the default in \texttt{code/calc-03.py}.
$R(t_{\mathrm{f}})$ is recorded for each $K$, and the first time $R(t)$ crosses $0.9$ is marked when that crossing occurs before $t_{\mathrm{f}}$ (\texttt{code/calc-03.py}).

\begin{figure*}[t]
    \centering
    \includegraphics[width=\textwidth]{figures/figure_3.pdf}
    \caption{Coupling-resolved synchronization diagnostics for the same homogeneous ODE setup as Figs.~\ref{fig:kkg-sync} and \ref{fig:kkg-order-K} ($H=0$, $\nabla\to 0$; no spatial boundaries), with $N=5$, $m_j\in[1.0,1.5]$, and the same fixed initial draw $\theta_j(0)\sim\mathcal{U}(-\pi,\pi)$, $\rho_j(0)\sim\mathcal{U}(0.8,1.2)$ for every $K$.
    (a) Terminal Kuramoto order parameter $R(t_{\mathrm{f}})$ with $t_{\mathrm{f}}=45$.
    (b) First crossing time $t_{0.9}$ with $R(t_{0.9})=0.9$ (linearly interpolated between time samples); curves omit $K$ for which the threshold is not reached before $t_{\mathrm{f}}$.
    Generated by \texttt{code/calc-03.py}.}
    \label{fig:kkg-K-scan}
\end{figure*}

The overdamped Kuramoto limit also highlights a competition between coupling and Hubble drag (App.~B), with dimensional criterion $K\gtrsim 3H^2N$ in Eq.~\eqref{eq:sync_threshold}.
Figure~\ref{fig:kkg-cosmo-sync} is not another homogeneous ODE integration; a spatially flat FRW background function $H(z)$ on $0\leq z\leq 45$ is used and compared to $K/(3H(z)^2N)$ from \texttt{code/calc-04.py} (defaults: $N=5$, $\Omega_m=0.315$, $\Omega_\Lambda=0.685$).

\begin{figure*}[t]
    \centering
    \includegraphics[width=\textwidth]{figures/figure_4.pdf}
    \caption{Schematic FRW background and cosmological synchronization proxy tied to Eq.~\eqref{eq:sync_threshold}.
    Spatially flat model with $\Omega_m=0.315$, $\Omega_\Lambda=0.685$, and $N=5$.
    (a) Hubble rate $H(z)/H_0$.
    (b) Ratio $K/\bigl(3H(z)^2N/H_0^2\bigr)$ for fixed illustrative $K/H_0$ on $0\leq z\leq 45$; the dotted line marks $K=3H(z)^2N$ in the same units (values $K/H_0=10,\,35,\,120$ from \texttt{code/calc-04.py} defaults).
    Curves above unity correspond to $K\gtrsim 3H(z)^2N$ at that redshift (Eq.~\eqref{eq:sync_threshold}).
    Generated by \texttt{code/calc-04.py}.}
    \label{fig:kkg-cosmo-sync}
\end{figure*}

Appendix~B also packages the linearized amplitude sector into an effective mass $m_{\mathrm{eff}}^2=m^2-K-\dot{\bar{\theta}}^2$ (Eq.~\eqref{eq:meff}) and states a Jeans-type instability when $m_{\mathrm{eff}}^2<0$ (Eq.~\eqref{eq:jeans_cond}).
Figure~\ref{fig:kkg-meff} uses a fixed $m=1$ normalization for axes; it is not generated from the ODE integrations of Figs.~\ref{fig:kkg-sync}--\ref{fig:kkg-K-scan}.
Defaults are $K\in[0,1.35]$, $\dot{\bar{\theta}}^2\in[0,m^2]$ on the map panel, and $\dot{\bar{\theta}}^2\in\{0,0.1,0.25\}$ in the slice curves (\texttt{code/calc-05.py}).

\begin{figure*}[t]
    \centering
    \includegraphics[width=\textwidth]{figures/figure_5.pdf}
    \caption{Schematic linear-stability diagnostics tied to Eqs.~\eqref{eq:meff} and \eqref{eq:jeans_cond}.
    (a) $m_{\mathrm{eff}}^2$ versus $K$ at fixed illustrative $\dot{\bar{\theta}}^2$; dashed horizontal line marks $m_{\mathrm{eff}}^2=0$, and dotted verticals mark $K=m^2-\dot{\bar{\theta}}^2$ for each slice when it lies in the plotted interval.
    (b) Colormap of $m_{\mathrm{eff}}^2$ in the $(K,\dot{\bar{\theta}}^2)$ plane with $m=1$ and black contour $m_{\mathrm{eff}}^2=0$.
    Generated by \texttt{code/calc-05.py}.}
    \label{fig:kkg-meff}
\end{figure*}

It is already shown in Eq.~\eqref{eq:14} that $K$ feeds the amplitude sector through $\cos(\theta_k-\theta_i)$ terms.
Figure~\ref{fig:kkg-stability} summarizes a stability-centric view in the same homogeneous ODE truncation ($H=0$, $\nabla\to 0$), using identical masses $m_j=m$ and the same initial draw and seed as Figs.~\ref{fig:kkg-sync}--\ref{fig:kkg-K-scan}, but with $m=1$ so that the tachyonic boundary lies at $K=m^2$ (Appendix~\ref{app:stability}).
The script is \texttt{code/calc-06.py}.

\begin{figure*}[!t]
    \centering
    \includegraphics[width=\textwidth]{figures/figure_6.pdf}
    \caption{Stability of the PLKG amplitude sector.
    (a) Mean amplitude in the stable regime $K<m^2$.
    (b) Effective mass squared from linear theory, Eq.~\eqref{eq:meff}, with $\dot{\bar{\theta}}=0$; the shaded region $K>m^2$ is tachyonic; the vertical dotted line marks $K=m^2$.
    (c) Regularization by a phenomenological cubic self-interaction $-\lambda(\rho_j-1)^3$ added to the toy amplitude equation: the $\lambda=0$ trajectory (dashed) exhibits the unphysical blow-up of the previous Fig.~6 at strong coupling, while $\lambda=0.1$ (solid) remains bounded.
    Simulation parameters: $N=5$, $m=1$, same seed as Figs.~\ref{fig:kkg-sync}--\ref{fig:kkg-K-scan}.
    Generated by \texttt{code/calc-06.py}.}
    \label{fig:kkg-stability}
\end{figure*}

The preceding Figs.~\ref{fig:kkg-sync}--\ref{fig:kkg-stability} suppress spatial gradients entirely.
To show that the same coupled amplitude--phase structure of Eqs.~\eqref{eq:14} and~\eqref{eq:16} can be stepped on a compact spatial domain, the Laplacian and gradients are discretized on a periodic cubic grid and $(\rho_i,\dot\rho_i,\theta_i,\dot\theta_i)$ are time-advanced with a fourth-order Runge--Kutta integrator in the local Minkowski patch ($a=1$, $H=0$).
In \texttt{code/calc-09.py}, $N$ fields are kept on $n_x^3$ cells (defaults $N=4$, $n_x=20$), two couplings $K=0$ and $K=3.5$ are compared with identical initial data and seed, and the interference density $|\Phi|^2$ with $\Phi=\sum_i \rho_i e^{i\theta_i}$ is visualized.
Optional Newtonian gravity is exposed only as a \emph{toy} channel for the figure pipeline: a fixed softened external potential, or a periodic FFT solve of Poisson's equation for a potential sourced by the total amplitude $\sum_i\rho_i$, each adding a phenomenological term $-\lambda\,\rho_i\,\Phi$ to the amplitude acceleration (not derived here from the covariant stress tensor).
The manuscript PDF uses the default \texttt{--gravity none}.
Figure~\ref{fig:kkg-3d-pde} reports the default run (\texttt{code/calc-09.py}).

\begin{figure*}[t]
    \centering
    \includegraphics[width=\textwidth]{figures/figure_9.pdf}
    \caption{Three-dimensional periodic-grid toy evolution of Eqs.~\eqref{eq:14} and~\eqref{eq:16} ($a=1$, $H=0$).
    (a)--(b) Mid-$z$ slices of $|\Phi|^2$ at the final time for $K=0$ and $K=3.5$ with the same initial data.
    (c) Peak value of $|\Phi|^2$ versus time for both couplings.
    (d) Text panel: variance of that peak time series after an early transient and a reminder of the optional gravity coupling implemented in \texttt{code/calc-09.py}.
    This panel is illustrative PDE numerics, not a fit to observations.}
    \label{fig:kkg-3d-pde}
\end{figure*}

\subsection{Behavior in Different Regimes}
\label{subsec:regimes}

\textbf{Strong coupling ($K \gg H$):}
Synchronization pressure dominates, suppressing structure
formation below $\lambda_{\text{sync}}$.

\textbf{Weak coupling ($K \ll H$):}
Model reduces to standard FDM with $w \approx 0$.

\subsection{Testable Predictions}
\label{subsec:predictions}

\begin{itemize}
    \item Halo mass function cutoff at $M_{\text{min}} \sim ...$
    \item Modified rotation curves in dwarf galaxies
    \item Lyman-$\alpha$ forest constraints: $K < ...$
\end{itemize}
